%!TEX program = uplatex
\RequirePackage{plautopatch}
\documentclass[uplatex,dvipdfmx]{jlreq}
\usepackage{amsmath,amssymb,amsfonts,amsthm}
\usepackage{enumerate}

\pagestyle{empty}

\begin{document}

%======================================================================
% 講演1：増田一男「正20面体とそのなかま」
%======================================================================
\begin{center}
{\large\bfseries 正20面体とそのなかま}\\[6pt]
\hspace*{\fill}{\large\bfseries 増田 一男（東工大院・理工）}
\end{center}

正20面体は正多面体の一つですが、正多面体の条件を緩めた色々な多面体が知
られています。それらについてお話をします。

まず正20面体の形はユニークでしょうか。
これについては Cauchy の剛体性定理があって凸多面体は鏡像を除いてユニーク
です。
凹でもユニークであろうと Euler は書いていますが、
1978年に Connelly は動く多面体を作りました。
動いている間体積は一定だろうという Sillvan の蛇腹予想は
1997年に Connelly らによって証明されました~[11]。

正多面体とは、(1) 凸で (2) 各面は合同な正多角形で
(3) 各頂点の回りは合同な正多角錐である多面体ですが、
5個あって Plato（プラトン）の立体と呼ばれています。
凸の条件をなくして星形正多角形や星形正多角錐を許すと
Kepler の小星形12面体と大星形12面体、
Poinsot（ポアンソ）の大12面体と大20面体の4個が加わります。
このうち2個は種数4の曲面です。
4個に限ることを Cauchy が彼の第2論文で示しています。

(1) 凸で (2') 各面は正多角形で
(3') 各頂点の回りは合同な多角錐である多面体を準正多面体といいます。
側面が正方形の正多角柱とそれをひねって側面を正3角形にした反正多角柱と
Platon の立体以外に $14+2$ 個あります。
(3') より強く (3'') 多面体をそれ自身に移す回転で任意の2頂点が移りあうとき
頂点推移的と言いますが、
Miller の多面体だけは頂点推移的ではありません。
これを除いた13個を Archimedes の立体といいます。

(2') 各面は（星形）正多角形で
(3'') 頂点推移的な多面体を一様多面体といい、
（星形）正多面体、準正多面体、星形（反）正多角柱以外に $54+7$ 個あります。
Coxeter（コクセター）らがこれで完全とは証明できていないが
これ以上はとめておけないと言って1954年に出したリスト~[7]に
1つ加えて1975年に Skilling~[8] がコンピューターを用いて完成しました。

(1) 凸で (2') 各面は正多角形である多面体は正多面体、
準正多面体以外に $91+5$ 個あります。
1966年に Johnson~[9] はリストを証明無しに発表し、
Zalgaller~[10] はコンピューターを用いて完全であることを示しました。

以上できちんとした多面体の分類は終わったように思えますが、
(1) 凸で (2'') 各面が合同な多角形である多面体は
まだ完全には分類されていないようです。
これは球面のタイル張りと考えられます。
準正多面体は外接球を持ち、
各頂点での接平面で囲まれた多面体はこの性質を持ちます。
Kepler の菱形12面体、菱形30面体はその例です。
(1)(2'') 多面体は2次元の moduli を持つことがあり、
実際合同な3角形20枚で20面体が作れます。

[1]はエピソードが多く楽しい本。文献豊富。[2]はすっきりとまとまった本。
[3]は不思議な絵本。6個の4次元正多面体の模型まで出ている。
[4] 紙模型の写真がすごい。[5]にはプログラムとグラフィックスが出ている。
[6]はがっちりした多面体の古典。

\begin{center}\textsc{参考文献}\end{center}
\begin{enumerate}[{[1]}]
  \item P.R.Cromwell, \textit{Polyhedra}, Cambridge, 1997.\\
    （邦訳）多面体, シュプリンガー・フェアラーク東京, 2001.
  \item 一松信, 正多面体を解く, 東海大学出版会, 2002.
  \item 宮崎興二, かたちと空間, 朝倉書店, 1983.
  \item M.J.Wenninger, \textit{Polyhedron Models}, Cambridge, 1971.\\
    （邦訳）多面体の模型, 教育出版, 1979.
  \item 関口次郎, 多面体の数理とグラフィックス, 牧野書店, 1996.
  \item H.S.M.Coxeter, \textit{Regular Polytopes}, Dover, 第3版1973.
  \item H.S.M.Coxeter, M.S.Longuet-Higgins, J.C.P.Miller,
    Uniform Polyhedra,
    \textit{Phil.\ Trans.\ Royal Soc.\ London}, Ser.~A \textbf{246} (1954) 401--450.
  \item J.Skilling, The complete set of uniform polyhedra,
    \textit{Phil.\ Trans.\ Royal Soc.\ London}, Ser.~A \textbf{278} (1975) 111--135.
  \item N.W.Johnson, Convex Polyhedra with Regular Faces,
    \textit{Can.\ J.\ M.} \textbf{18} (1966) 169--200.
  \item V.A.Zalgaller, Convex polyhedra with regular faces (in Russian),
    \textit{Seminars in Math., V.A.Steklov Math.\ Institute, Leningrad}, vol.~2, 1966.
    （英訳）Consultants Bureau, 1969.
  \item R.Connelly, I.Sabitov, A.Walz, The Bellows Conjecture,
    \textit{Contributions to Algebra and Geometry} \textbf{38} no.~1 (1997) 1--10.
\end{enumerate}

\newpage

%======================================================================
% 講演2：加藤文元「正20面体とリーマン面」
%======================================================================
\begin{center}
{\large\bfseries 正20面体とリーマン面}\\[6pt]
\hspace*{\fill}{\large\bfseries 加藤 文元（京大・理）}
\end{center}

リーマン球面の正則同型群、つまり一次分数変換群の中には5次交代群
と同型な部分群が共役を除いてただ一つあります。5次交代群の各元の位数は
2, 3, 5ですが、これらの元の固定点をリーマン球面上にプロットしていくと、
正20面体の姿が浮き上がります。具体的には位数5の元の固定点が丁度正20面
体の頂点になっています。さてリーマン球面をこの5次交代群で割ると、また
リーマン球面になるのですが、このとき得られるリーマン球面からリーマン球
面への写像は、3点で分岐指数が 2, 3, 5 で分岐する被覆写像になってます。この
写像やその逆写像（多価）は函数としても面白いものです。ガウスの超幾何微分
方程式と深い関係にあります。以上のことは、つまり5次交代群がいわゆるシュ
ワルツの三角群というものの一つの例になっていることを示しています。もっ
と一般の三角群を考えると、色々面白い函数が得られます。例えば典型的な例
はラムダ函数とか $j$ 函数などです。以上のことについてお話しします。

\vfill

\newpage

%======================================================================
% 講演3：加藤文元「$p$-進20面体」
%======================================================================
\begin{center}
{\large\bfseries $p$-進20面体}\\[6pt]
\hspace*{\fill}{\large\bfseries 加藤 文元（京大・理）}
\end{center}

5次交代群のリーマン球面への作用の固定点を結ぶ、ということで正20
面体が出来たのですが、この操作を $p$-進幾何（あるいはリジッド幾何）という世
界で行うことが出来ます。$p$-進幾何におけるリーマン球面の類似としては、
Bruhat--Tits 樹木というものを考えます。これによって $p$-進20面体がある種の
樹木という形で作られるわけです。このことは、もう少し一般的な三角群につ
いても行うことが出来ます。$p$-進幾何の考え方から Bruhat--Tits 樹木が出てく
る理由も含めて、以上のことについてお話ししようと思います。

\vfill

\begin{center}\textsc{参考文献}\end{center}
\begin{enumerate}[{[1]}]
  \item G.~Cornelissen, F.~Kato,
    $p$-adic icosahedron,
    \textit{Notices AMS} \textbf{52}, No.~7, 720--727 (2005).
  \item F.~Kato,
    Non-archimedean orbifolds covered by Mumford curves,
    \textit{Journal of Alg.\ Geom.} \textbf{14} (2005) 1--34.
  \item 加藤文元, RIGID幾何と一意化、2003年日本数学会（千葉大学）代数分科会講演アブストラクト、特別講演4.
\end{enumerate}

\vfill

\newpage

%======================================================================
% 講演4：橋本義武「クライン「正20面体と5次方程式」を読む」
%======================================================================
\begin{center}
{\large\bfseries クライン「正20面体と5次方程式」を読む}\\[6pt]
\hspace*{\fill}{\large\bfseries 橋本 義武（阪市大・理）}
\end{center}

一般の5次以上の代数方程式が根号によって解けないことは
よく知られていますが、
ここでは根号以外の操作を許して5次方程式を解くことを考えます。
それには、正二十面体に関連するある特殊関数を用います。

また、有限体の元を成分とする2行2列の行列に関するガロワの定理と、
正多面体との関係についても少し触れたいと思います。

\begin{center}\textsc{参考文献}\end{center}
\begin{enumerate}[{[1]}]
  \item F.~クライン, 正20面体と5次方程式, シュプリンガー・フェアラーク東京, 2005.
  \item J.-P.~Serre, Extensions icosaédriques, 1980, in \textit{Collected Papers}, vol.~III.
  \item 鈴木通夫, 群論（上）, 岩波書店, 1977.
\end{enumerate}

\vfill

\newpage

%======================================================================
% 講演5：橋本義武「グロタンディークの双二十面体――マチウ群試論」
%======================================================================
\begin{center}
{\large\bfseries グロタンディークの双二十面体――マチウ群試論}\\[6pt]
\hspace*{\fill}{\large\bfseries 橋本 義武（阪市大・理）}
\end{center}

対称群のうち6次のものだけは、外部自己同型をもちます。
その機構を説明するのが、
グロタンディークによって定義された双二十面体です。
6点集合上には6つの双二十面体構造が入り、
これらが「双対」6点集合をなします。
6次対称群が外部自己同型をもつ理由は、
「6点集合は『双対』をもつ」という事実にあるのです。

6点集合とその「双対」の和である12点集合には
「シンプレクティック構造」が入ります。
シンプレクティック12点集合の自己同型群が、
散在型単純群の一つである12次マチウ群です。

さらにシンプレクティック12点集合は
「双対」シンプレクティック12点集合をもち、
それらの和である24点集合にもある構造が誘導されます。
その自己同型群が24次マチウ群です。

\begin{center}\textsc{参考文献}\end{center}
\begin{enumerate}[{[1]}]
  \item A.~グロタンディーク, ある夢と数学の埋葬, 現代数学社, 1993.
  \item 大山豪, 有限置換群, 裳華房, 1981.
  \item T.~Beth, D.~Jungnickel, H.~Lenz,
    \textit{Design Theory}, vol.~1, 2nd ed., Cambridge, 1999.
  \item R.~L.~Griess, Jr.,
    \textit{Twelve Sporadic Groups}, Springer-Verlag, 1998.
\end{enumerate}

\vfill

\end{document}